% Sección 5
\section{Lo que dicen los actores}

Las fuentes públicas nombran a pocos actores con nombre propio. En la Alcaldía de Manizales hablaron el alcalde, Jorge Eduardo Rojas, y la gestora social, Juliana Londoño, que coordinó el acopio del Coliseo Menor. En la Gobernación de Caldas aparece el gobernador, Henry Gutiérrez. En ABACO hablaron el director ejecutivo, Juan Carlos Buitrago, y la gerente de Cooperación Internacional, Norma Fernanda Alonso. El censo de damnificados lo coordinaron Bomberos, Defensa Civil y Cruz Roja con la Alcaldía, y lo acompañaron la Personería, la Defensoría, el ICBF y la fuerza pública.

Otros actores clave no tienen declaraciones públicas que hayamos encontrado. Es el caso del Banco de Alimentos de Manizales, que la investigación asocia con la Arquidiócesis sin haberlo confirmado. Tampoco encontramos declaraciones de la Cruz Roja de Caldas, la Defensa Civil de Caldas, la UDEGER ni la UGR de Manizales. La búsqueda de lo que dijeron los actores en redes, prensa e informes de gestión está \eninvestigacion.

La tabla~\ref{tab:entrevista} recoge las ocho preguntas que el equipo quiere hacer en entrevistas a quienes despacharon, pidieron o transportaron comida después del sismo. Ninguna respuesta viene todavía de una entrevista, porque el equipo no ha hablado con el Banco de Alimentos de Manizales ni con los organismos de socorro de Caldas. Lo que sabemos sale de prensa e informes públicos.

{\footnotesize
\setlength{\tabcolsep}{4pt}
\renewcommand{\arraystretch}{1.3}
\begin{longtable}{L{0.025\linewidth} L{0.17\linewidth} L{0.50\linewidth} L{0.22\linewidth}}
\caption{Las ocho preguntas de entrevista, con lo que ya sabemos y su fuente.}\label{tab:entrevista}\\
\toprule
\textbf{\#} & \textbf{Pregunta} & \textbf{Lo que ya sabemos} & \textbf{Fuente}\\
\midrule
\endfirsthead
\multicolumn{4}{l}{\small\itshape Tabla~\ref{tab:entrevista}, continuación}\\
\toprule
\textbf{\#} & \textbf{Pregunta} & \textbf{Lo que ya sabemos} & \textbf{Fuente}\\
\midrule
\endhead
\midrule
\multicolumn{4}{r}{\small\itshape Sigue en la página siguiente}\\
\endfoot
\bottomrule
\endlastfoot
1 & ¿Cómo decidían el despacho? & ABACO anunció que daría prioridad a los albergues oficiales, los hospitales y las familias autoalbergadas. Según la investigación de referentes, ABACO registra fechas de vencimiento; la fuente directa de ese dato está \porverificar. Cómo decide paso a paso el Banco de Alimentos de Manizales está \pendiente. & Noticias RCN, 10 y 11 ago 2026 (director ejecutivo de ABACO)\\
2 & ¿Cómo llegaban los pedidos? & \pendiente. No hay información pública sobre cómo pedían los comités barriales en Manizales. Que los pedidos lleguen desordenados por WhatsApp y por teléfono es una hipótesis del equipo. Como referencia, en Rio Grande do Sul (Brasil, 2024) los pedidos de rescate circularon por grupos de WhatsApp y se repetían. & Fast Company Brasil (referente regional)\\
3 & ¿Cómo se enteraban de los cierres? & La prensa citó reportes de la Policía de Tránsito, y la ANI informó el estado de sus corredores. No encontramos comunicados directos del INVÍAS, la Gobernación ni la Alcaldía. Que el reporte del conductor viaje por teléfono, tarde y sin registro es una hipótesis del equipo. Cómo se enteraban quienes despachaban está \pendiente. & Publimetro, Semana y Colombia.com (10 y 12 ago 2026); ANI (10 ago 2026)\\
4 & ¿Qué comida se dañó? & ABACO explicó que sin energía no podía manejar perecederos y que priorizó enlatados y alimentos de larga vida. La Cruz Roja pidió no donar perecederos. La cifra de alimentos donados que se dañaron o vencieron en Manizales está \pendiente. & La República (entrevista, sin fecha en la investigación); Infobae, 14 ago 2026\\
5 & ¿Qué plataformas usaron? & En el país circularon Cuidar a Colombia (EAFIT), un bot de WhatsApp de JumpCube con AWS y los botones de donación de ABACO y Rappi. El Estado usó el Registro Único de Damnificados, que cerró el 4 de septiembre. Cuáles usaron los actores de Manizales está \pendiente. & El Colombiano; Noticias RCN; Infobae, 11 ago 2026; Vanguardia, 5 oct 2026\\
6 & ¿Qué les haría falta? & ABACO nombró la llegada al territorio, el transporte disponible, la articulación entre actores, el almacenamiento en los acopios y vías de acceso claras hacia las zonas rurales. El acopio de El Campín pidió camiones, montacargas y voluntarios. Con corte al 13 de agosto, Manizales pedía colchonetas, mantas, carpas y plantas eléctricas para el albergue. Lo que diría el Banco de Alimentos de Manizales está \pendiente. & Noticias RCN, 18 ago 2026; Noticias RCN (El Campín, sin fecha); Asocapitales n.º 24\\
7 & ¿Quién usaría la herramienta? & La ficha del reto nombra a la Cruz Roja, la Defensa Civil, el Banco de Alimentos, los transportadores y los comités barriales. El equipo propone a ABACO y al Banco de Alimentos como usuario principal. La confirmación de los propios actores está \pendiente. & Ficha de la vertical 04, citada en \ruta{docs/PROBLEMA.md}; \ruta{docs/ESTADO_DEL_ARTE.md}\\
8 & ¿Quién nos recibiría para un piloto? & Los candidatos son ABACO, el Banco de Alimentos de Manizales, la UGR de Manizales, la Cruz Roja de Caldas, el PMA, el equipo humanitario de OpenStreetMap (HOT), el IDEA de la UNAL y Asocapitales. El nombre y el canal oficial del Banco de Alimentos de Manizales no están confirmados. Ninguna respuesta todavía, así que está \pendiente. & \ruta{docs/investigacion/05_datos_abiertos_y_aliados.md}, parte C\\
\end{longtable}
}
